% Class-compatibility preamble retained for a reader excerpt of
%   Lisa Sauermann, Dmitrii Zakharov, "A Sharp Ramsey Theorem for Ordered
%   Hypergraph Matchings", arXiv:2309.04813v4, CC BY 4.0.
%
% The source pins the publisher class aic (the AIMS journal class), which is not part of
% the delivery baseline, so this file supplies the amsart baseline and then copies the
% author's own declaration block mechanically and unchanged
% (pinned-original-source.tex lines 25 and 28-59: the ams packages, the author's eleven
% newtheorem declarations with [theorem] as the shared counter, the definition theorem
% style, the author's comment that the file should add as few macro or package
% definitions as possible, and the author's own macros \su, \sm, \Z, \HH, \conv and \PP).
% The author's tikz package (pinned source line 26) serves the three figures of the paper
% that this excerpt does not carry and is therefore not reproduced; the author's babel,
% fontenc and inputenc lines (pinned source lines 47-49) are layout/font choices of the
% publisher build and are not part of the mathematical excerpt, the delivery engine being
% XeLaTeX with UTF-8 input.
% No mathematical macro, statement, symbol or assumption is introduced, changed or
% removed by this file.
\documentclass[10pt,a4paper]{amsart}
\usepackage{amsthm,amsfonts,amssymb,amsmath}
\usepackage{graphicx}

\newtheorem{theorem}{Theorem}[section]
\newtheorem{proposition}[theorem]{Proposition}
\newtheorem{lemma}[theorem]{Lemma}
\newtheorem{conjecture}[theorem]{Conjecture}
\newtheorem{fact}[theorem]{Fact}
\newtheorem{claim}[theorem]{Claim}
\newtheorem{corollary}[theorem]{Corollary}
\newtheorem{definition}[theorem]{Definition}
\newtheorem{problem}[theorem]{Problem}
\newtheorem{question}[theorem]{Question}
\newtheorem{observation}[theorem]{Observation}

\theoremstyle{definition}
\newtheorem{remark}[theorem]{Remark}

% The author notes at this point that as few macro or package
% definitions as possible should be added; the six macros below are the
% author's own and are copied unchanged.

\newcommand{\su}{\subseteq}
\newcommand{\sm}{\setminus}
\newcommand{\Z}{\mathbb{Z}}
\newcommand{\HH}{\mathcal{H}}
\newcommand{\conv}{\operatorname{conv}}
\newcommand{\PP}{\mathcal{P}}
\usepackage{cleveref}
\crefname{originaldefinition}{original definition}{original definitions}
\Crefname{originaldefinition}{Original definition}{Original definitions}
\makeatletter
\newlabel{def-s-partite}{{2.1}{source}{}{sourcepointer-1}{}}
\newlabel{def-s-partite@cref}{{[originaldefinition][0][]2.1}{[1][0][]source}{}{}{}}
\makeatother
\crefname{originallemma}{original lemma}{original lemmas}
\Crefname{originallemma}{Original lemma}{Original lemmas}
\makeatletter
\newlabel{obs-interval-wise}{{3.1}{source}{}{sourcepointer-2}{}}
\newlabel{obs-interval-wise@cref}{{[originallemma][0][]3.1}{[1][0][]source}{}{}{}}
\makeatother
\crefname{originaltheorem}{original theorem}{original theorems}
\Crefname{originaltheorem}{Original theorem}{Original theorems}
\makeatletter
\newlabel{thm-paramteters-upper-bound}{{1.2}{source}{}{sourcepointer-3}{}}
\newlabel{thm-paramteters-upper-bound@cref}{{[originaltheorem][0][]1.2}{[1][0][]source}{}{}{}}
\makeatother
\title{A dichotomy for dense lambda-partite ordered hypergraph collections with bounded clique sizes}
\author{Lisa Sauermann, Dmitrii Zakharov}
\date{arxiv:2309.04813v4; CC BY 4.0}
\begin{document}
\maketitle
\noindent\textit{Editorial note.} All delivered mathematical text is the author's own, copied from the pinned author source sources/combinatorics-fl02-a-sharp-ramsey-theorem-for-ordered/pinned-original-source.tex; no statement or proof step was rewritten, reordered or removed.\par\medskip
\noindent\textit{Editorial note.} Preamble provenance: the source pins the publisher class aic, which is not part of the delivery baseline, so this card uses sources/combinatorics-fl02-a-sharp-ramsey-theorem-for-ordered/retained-preamble.tex, which supplies the amsart baseline and then copies the author's ams packages, the author's eleven newtheorem declarations sharing the theorem counter, the definition theorem style and the author's own macros \textbackslash{}su, \textbackslash{}sm, \textbackslash{}Z, \textbackslash{}HH, \textbackslash{}conv and \textbackslash{}PP (pinned source lines 25 and 28-59) verbatim. No macro was invented and no mathematical symbol was changed.\par\medskip
\noindent\textit{Editorial note.} Bibliography provenance: the citations of the counted proof are resolved with sources/combinatorics-fl02-a-sharp-ramsey-theorem-for-ordered/bibliography.tex, an excerpt of the author's own thebibliography environment (pinned source lines 598-703) copied character for character with all eighteen entries kept; no entry was rewritten, reordered, merged, split or invented and no citation key was changed.\par\medskip
\noindent\textit{Editorial note.} External references are navigation pointers only; every deliberately kept step of the author's argument stays inside the counted spans. Their printed numbers are the author's own numbering under the author's newtheorem declarations, computed mechanically by the lane's scanner.\par\medskip
\noindent\textit{Editorial note.} The labels the counted proof invokes outside its own span, def-s-partite, obs-interval-wise and thm-paramteters-upper-bound, are navigation-pointer stubs; obs-interval-wise is also counted by another line's card on the other staged copy of this paper, which is why it is a pointer here rather than a second counted claim.\par\medskip
\noindent\textit{Allowed standard prerequisites (editorial).}\begin{itemize}
\item Ordered r-patterns, the pattern family of an ordered partition, P-cliques and their sizes, and the notion of a lambda-partite collection. Set up in the author's preliminaries (source lines 225-281) and in the author's definition at source lines 289-291. Fixes every object of the statement: the collection, its pattern family and the clique-size bounds.
\end{itemize}\medskip
\ifdefined\hypertarget\hypertarget{sourcepointer-1}{}\fi
\noindent\textit{Source reference.} Original definition 2.1: sources/combinatorics-fl02-a-sharp-ramsey-theorem-for-ordered/pinned-original-source.tex, source lines 289--289. Navigation pointer only: the counted proof refers to the author's definition of a partite collection; its content is declared as a named standard prerequisite and is not reproduced.\par\medskip
\ifdefined\hypertarget\hypertarget{sourcepointer-2}{}\fi
\noindent\textit{Source reference.} Original lemma 3.1: sources/combinatorics-fl02-a-sharp-ramsey-theorem-for-ordered/pinned-original-source.tex, source lines 359--359. Navigation pointer only: the counted proof invokes the author's interval-wise lemma; that statement is counted by another line's card LIT-comb-omr-interval-wise on the other staged copy of this paper, so it is a pointer here and is declared as a named standard prerequisite.\par\medskip
\ifdefined\hypertarget\hypertarget{sourcepointer-3}{}\fi
\noindent\textit{Source reference.} Original theorem 1.2: sources/combinatorics-fl02-a-sharp-ramsey-theorem-for-ordered/pinned-original-source.tex, source lines 199--199. Navigation pointer only: the counted proof refers to the author's upper-bound theorem for the sizes appearing in the alternatives; it is not reproduced.\par\medskip
\begin{proposition}\label{prop-two-options}
Let $r$ be a positive integer, and let $\lambda=(\lambda_1,\dots,\lambda_s)$ be an ordered partition of $r$ into $s<r$ parts. For every ordered partition $\lambda'$ of $r$ into $s+1$ parts with $\lambda\succ\lambda'$, consider some positive integer $M_{\lambda'}$, and let $M=\sum_{\lambda'} M_{\lambda'}$ be the sum of all these positive integers. Let $\HH$ be a finite collection of pairwise disjoint subsets $e\su \Z$, each of size $|e|=r$, such that $\HH$ is $\lambda$-partite. Then at least one of the following two statements holds:
\begin{itemize}
    \item[(a)] There is a sub-collection $\HH^*\su \HH$ of size $|\HH^*|\ge |\HH|/(2M)$ such that $\HH^*$ is interval-wise $\lambda$-partite.
    \item[(b)] For some ordered partition $\lambda'$ of $r$ into $s+1$ parts with $\lambda\succ\lambda'$, there is a sub-collection $\HH'\su \HH$ of size $|\HH'|>M_{\lambda'}$ such that $\HH'$ is $\lambda'$-partite.
\end{itemize}
\end{proposition}

\begin{proof}[Proof of Proposition~\ref{prop-two-options}]
Suppose that $\HH$ does not have a subcollection $\HH'$ as in (b). We need to show that then (a) holds. Let $x_0<x_1<\dots<x_s$ be as in Definition~\ref{def-s-partite}. For every $e\in \HH$, we consider the intervals $e\cap (x_{i-1},x_i)$ for $i=1,\dots,s$. Our goal is to find a subcollection $\HH^*\su \HH$ of size $|\HH^*|\ge |\HH|/(2M)$ such that for each $i=1,\dots,s$, the intervals $\conv(e\cap (x_{i-1},x_i))$ are pairwise disjoint for all $e\in \HH^*$. Note that for any $i=1,\dots,s$ and any $e\in \HH$, the interval $\conv(e\cap (x_{i-1},x_i))$ is a closed sub-interval of $(x_{i-1},x_i)$. Furthermore, the endpoints of all these sub-intervals are distinct (as all $e\in \HH$ are pairwise disjoint subsets of $\Z$).

Let us consider an auxiliary graph $G$ with vertex set $\HH$, where for $e,e'\in\HH$ we draw an edge between $e$ and $e'$ if $\conv(e\cap (x_{i-1},x_i))\cap \conv(e'\cap (x_{i-1},x_i))\ne \emptyset$ for some $i\in\{1,\dots,s\}$. Our goal is then to show that the graph $G$ has an independent set of size at least $|\HH|/(2M)$.

Note that for any two closed intervals $[a,b]$ and $[c,d]$ with $[a,b]\cap [c,d]\ne \emptyset$ we have $a\in [c,d]$ or $c\in [a,b]$. Hence we can orient the edges of $G$ in such a way that for any $e,e'\in\HH$ with an edge from $e$ to $e'$ there is an index $i\in\{1,\dots,s\}$ with $\min(e\cap (x_{i-1},x_i))\in \conv(e'\cap (x_{i-1},x_i))$.

In order to show that the graph $G$ contains an independent set of size at least $|\HH|/(2M)$, it suffices to show that $G$ has average degree at most $2M-1$ (by a well-known result of Caro~\cite{Caro} and Wei~\cite{Wei}). Considering the orientation of the edges of $G$, it therefore suffices to show that for every $e\in \HH$ there are at most $M-1$ outgoing edges at $e$ in the graph $G$ (then there are at most $(M-1)|\HH|$ total edges in $G$ and hence the average degree is at most $2M-2$).

So let us now fix some $e\in \HH$. We need to show that there are at most $M-1$ different $e'\in\HH\sm\{e\}$ such that $\min(e\cap (x_{i-1},x_i))\in \conv(e'\cap (x_{i-1},x_i))$ for some index $i\in\{1,\dots,s\}$. 

For every index $i\in\{1,\dots,s\}$ with $\lambda_i=1$, we have $|e'\cap (x_{i-1},x_i)|=\lambda_i=1$ for every $e'\in \HH$, and the singleton sets $e'\cap (x_{i-1},x_i)$ are all distinct for all $e'\in \HH$ (as all $e'\in \HH$ are pairwise disjoint). Hence we cannot have $\min(e\cap (x_{i-1},x_i))\in \conv(e'\cap (x_{i-1},x_i))$ for any $e'\in \HH\sm\{e\}$ when $\lambda_i=1$.

We claim that for each $i\in\{1,\dots,s\}$ with $\lambda_i\ge 2$, there can be at most $\sum_{\lambda'} M_{\lambda'}-1$ different $e'\in\HH\sm\{e\}$ with $\min(e\cap (x_{i-1},x_i))\in \conv(e'\cap (x_{i-1},x_i))$, where the sum is over all ordered partitions $\lambda'$ of $r$ into $s+1$ parts that can be obtained from $\lambda$ by dividing the $i$-th part $\lambda_i$ into two pieces. This suffices to finish the proof, since then the total number of $e'\in\HH\sm\{e\}$ with $\min(e\cap (x_{i-1},x_i))\in \conv(e'\cap (x_{i-1},x_i))$ for some $i\in\{1,\dots,s\}$ is at most
\[\sum_{\substack{i\in\{1,\dots,s\}\\\lambda_i\ge 2}}\left(\left(\sum_{\substack{\lambda'\text{obtained from }\lambda\text{ by}\\\text{dividing }\lambda_i\text{into two parts}}}\!\!M_{\lambda'}\right)-1\right)=\sum_{\lambda'} M_{\lambda'}-|\{i\in\{1,\dots,s\}\mid \lambda_i\ge 2\}|\le \sum_{\lambda'} M_{\lambda'}-1=M-1,\]
where in the second and third term the sum is over all ordered partitions $\lambda'$ of $r$ into $s+1$ parts with $\lambda\succ\lambda'$ (note that each of these ordered partitions is obtained by dividing some part $\lambda_i$ of $\lambda$ with $\lambda_i\ge 2$ into two parts). Also note that for the inequality sign we have used that there must at least one index $i\in \{1,\dots,s\}$ with $\lambda_i\ge 2$, since $s<r$.

It remains to prove the above claim, so fix some index $i\in\{1,\dots,s\}$ with $\lambda_i\ge 2$. There are $\lambda_i-1$ different ordered partitions $\lambda'$ that can be obtained from $\lambda$ by dividing $\lambda_i$ into two parts, namely the ordered partitons $\mu^{(1)}=(\lambda_1,\dots,\lambda_{i-1},1,\lambda_i-1,\lambda_{i+1},\dots,\lambda_s), \ \mu^{(2)}=(\lambda_1,\dots,\lambda_{i-1},2,\lambda_i-2,\lambda_{i+1},\dots,\lambda_s),\ \dots,\ \linebreak \mu^{(\lambda_i-1)}=(\lambda_1,\dots,\lambda_{i-1},\lambda_i-1,1,\lambda_{i+1},\dots,\lambda_s)$. Defining $y=\min(e\cap (x_{i-1},x_i))\in \Z$, we need to show that there are at most $\sum_{j=1}^{\lambda_i-1} M_{\mu^{(j)}}-1$ different $e'\in\HH\sm\{e\}$ with $y\in \conv(e'\cap (x_{i-1},x_i))$. Note that for every such $e'\in\HH\sm\{e\}$, we also have $y+0.1\in \conv(e'\cap (x_{i-1},x_i))$, since $\conv(e'\cap (x_{i-1},x_i))$ is a closed interval with endpoints in $\Z\sm\{y\}$. We furthermore also have $y+0.1\in \conv(e\cap (x_{i-1},x_i))$ (recalling that $y=\min(e\cap (x_{i-1},x_i))$ and that $e\cap (x_{i-1},x_i)\su \Z$ is a set of size $\lambda_i\ge 2$), which in particular implies $x_{i-1}<y+0.1<x_i$. It now suffices to show that there are at most $\sum_{j=1}^{\lambda_i-1} M_{\mu^{(j)}}$ different $e'\in\HH$ with $y+0.1\in \conv(e'\cap (x_{i-1},x_i))$. For every such $e'$, we have $1\le |e'\cap (x_{i-1},y+0.1)|\le \lambda_i-1$ (recalling that $e'\su \Z$ and $|e'\cap (x_{i-1},x_i)|=\lambda_i$).

For each $j=1,\dots,\lambda_i-1$, there can be at most $M_{\mu^{(j)}}$ different $e'\in\HH$ with $|e'\cap (x_{i-1},y+0.1)|=j$. Indeed, for each such $e'\in\HH$ we have $|e'\cap (y+0.1,x_i)|=|e'\cap (x_{i-1},x_i)|-|e'\cap (x_{i-1},y+0.1)|=\lambda_i-j$. Hence, by considering $x_0,\dots,x_{i-1}, y+0.1,x_i,\dots,x_s\in \mathbb{R}\sm \Z$, we can observe that the collection $\HH'\su \HH$ of all $e'\in\HH$ with $|e'\cap (x_{i-1},y+0.1)|=j$ is $\mu^{(j)}$-partite. So by our assumption that (b) does not hold, there can be at most $M_{\mu^{(j)}}$ different such $e'\in\HH$ (recall that $\mu^{(j)}$ is an ordered partition of $r$ into $s+1$ parts with $\lambda\succ \mu^{(j)}$).

Thus, in total there can be at most $\sum_{j=1}^{\lambda_i-1} M_{\mu^{(j)}}$ different $e'\in\HH$ with $1\le |e'\cap (x_{i-1},y+0.1)|\le \lambda_i-1$, meaning that there are at most $\sum_{j=1}^{\lambda_i-1} M_{\mu^{(j)}}$ different $e'\in\HH$ with $y+0.1\in \conv(e'\cap (x_{i-1},x_i))$. This finishes the proof of the proposition.
\end{proof}

% Bibliography excerpt retained for reader excerpts of
%   Lisa Sauermann, Dmitrii Zakharov, "A Sharp Ramsey Theorem for Ordered Hypergraph Matchings",
%   arXiv:2309.04813v4, CC BY 4.0.
%
% Every entry below is copied character for character from the author's own thebibliography
% environment of the pinned author source (lines 598-703).  No entry was rewritten, reordered,
% merged, split or invented and no citation key was changed; all 18 entries are kept so that
% any citation of a delivered excerpt resolves without guessing an entry.
\begin{thebibliography}{99}

%\bibitem{bergelson-johnson-moreira}
%Vitaly Bergelson, John H. Johnson Jr., and Joel Moreira.
%\newblock New polynomial and multidimensional extensions of classical partition
%  results.
%\newblock 2015, arXiv:1501.02408.

%\bibitem{cilleruelo}
%Javier Cilleruelo.
%\newblock Combinatorial problems in finite fields and {S}idon sets.
%\newblock {\em Combinatorica}, 32(5):497--511, 2012.


\bibitem{AJKS} 
Michael Anastos, Zhihan Jin, Matthew Kwan, and Benny Sudakov.
\newblock Extremal, enumerative and probabilistic results on ordered hypergraph matchings.
\newblock {\em Forum of Mathematics Sigma} 13: Paper No.~e55, 39~pp, 2025.

\bibitem{Campos} 
Marcelo Campos, Simon Griffiths, Robert Morris, and Julian Sahasrabudhe. 
\newblock An exponential improvement for diagonal Ramsey.
\newblock 2023, arXiv:2303.09521.

\bibitem{Caro} 
Yair Caro. 
\newblock New Results on the Independence Number.
\newblock Technical Report, Tel-Aviv University, 1979.

\bibitem{ExpES} 
Recep Altar Çiçeksiz, Recep Altar, Zhihan Jin, Eero Räty, and István Tomon. 
\newblock Exponential Erd\H{o}s-Szekeres theorem for matrices.
\newblock {\em Proceedings of the London Mathematical Society} 129, no. 3: Paper No. e12632, 14~pp, 2024.

\bibitem{Dil} 
Robert P. Dilworth 
\newblock A Decomposition Theorem for Partially Ordered Sets.
\newblock {\em Annals of Mathematics}  51: 161--166, 1950.

\bibitem{DGR22}
Andrzej Dudek, Jarosław Grytczuk, and Andrzej Ruci\'{n}ski. 
\newblock Ordered unavoidable sub-structures in matchings and random matchings.
\newblock {\em Electronic Journal of Combinatorics} 31, no. 2: Paper No. 2.15, 27 pp, 2024

\bibitem{DGR} 
Andrzej Dudek, Jarosław Grytczuk, and Andrzej Ruci\'{n}ski. 
\newblock Erd\H{o}s-Szekeres type Theorems for ordered uniform matchings.
\newblock {\em Journal of Combinatorial Theory, Series B} 170: 225--259, 2025.

\bibitem{ES} 
Paul Erd{\H o}s and George Szekeres. 
\newblock A combinatorial problem in geometry.
\newblock {\em Compositio Mathematica} 2: 463--470, 1935.

\bibitem{GL} Timothy W. Gowers and Jason Long.
\newblock The length of an s-increasing sequence of r-tuples. \newblock {\em Combinatorics, Probability and Computing} 30, no. 5: 686--721, 2021.

\bibitem{graham-rothschild-spencer} 
Ronald L. Graham, Bruce L. Rothschild, and Joel H. Spencer. 
\newblock {\em Ramsey theory.} Vol. 20. John Wiley \& Sons, 1991.


\bibitem{Lee} 
Choongbum Lee. 
\newblock Ramsey numbers of degenerate graphs.
\newblock {\em Annals of Mathematics} 185, no. 3: 791--829, 2017.

\bibitem{LS} 
Nathan Linial and Michael Simkin. 
\newblock Monotone subsequences in high-dimensional permutations. 
\newblock {\em Combinatorics, Probability and Computing} 27, no. 1: 69--83, 2018.

\bibitem{MV} 
Sam Mattheus and Jacques Verstraete. 
\newblock The asymptotics of $r(4,t)$.
\newblock {\em Annals of Mathematics} 199, no.~2: 919–941,  2024.

%\bibitem{Mirsky} 
%Leon Mirsky. 
%\newblock A dual of Dilworth’s decomposition theorem.
%\newblock {\em Amer. Math. Monthly} 78, no. 8: 876-877, 1971.

\bibitem{MS} 
Guy Moshkovitz and Asaf Shapira. 
\newblock Ramsey theory, integer partitions and a new proof of the Erdős–Szekeres theorem. 
\newblock {\em Advances in Mathematics} 262: 1107--1129, 2014.

\bibitem{promel} 
Hans Jürgen Prömel. 
\newblock {\em Ramsey theory for discrete structures.} Vol. 3. New York: Springer, 2013.

\bibitem{RR} 
Christian Reiher, and Vojt\v{e}ch R\"odl. 
\newblock The girth Ramsey theorem.
\newblock 2023, arXiv:2308.15589.

\bibitem{ST} 
Tibor Szab\'{o} and G\'{a}bor Tardos. 
\newblock A multidimensional generalization of the Erd\H{o}s–Szekeres lemma on monotone subsequences.
\newblock {\em Combinatorics, Probability and Computing} 10, no. 6: 557--565, 2001.

\bibitem{Wei} 
Victor K. Wei. 
\newblock A Lower Bound on the Stability Number of a Simple Graph.
\newblock {\em Technical memorandum}, TM 81-11217-9, Bell Laboratories, 1981.


\end{thebibliography}
\end{document}
